% answers_short.tex —— 2014 年·湖南卷·文科 · 简短答案速查
% 由 tools/make_short_answers.py 自 answers.tex 的 \gksolhead/\gkanswer 提取，
% 并与 2014/2014_数学真题_湖南文_答案要点.md 逐题比对后生成，勿手改（改源 MD 后重跑）。

\begin{gknotice}
\textbf{使用说明}：本节只给最终答案便于快速核对，完整推导见解析版。答案为编写者独立推导，\textbf{非官方标准答案}。
\end{gknotice}

\gkpart{选择题}
\begin{center}\begin{tabular}{ccccccccccc}
\toprule
题号 & 1 & 2 & 3 & 4 & 5 & 6 & 7 & 8 & 9 & 10 \\
\midrule
答案 & B & C & D & A & B & C & D & B & C & D \\
\bottomrule
\end{tabular}\end{center}

\gkpart{填空题}
\begin{itemize}[leftmargin=2.6em]
  \item[\textbf{11.}] $-3$。$\mathrm i^2=-1$，故 $\dfrac{3+\mathrm i}{\mathrm i^2}=-3-\mathrm i$，实部为 $-3$。
  \item[\textbf{12.}] $x-y-1=0$。两式相减消去 $t$ 得 $x-y=1$。
  \item[\textbf{13.}] $7$。可行域的三个顶点为 $(1,1)$，$(3,1)$，$(2,2)$，代入 $z=2x+y$ 分别得 $3$，$7$，$6$，最大值为 $7$。
  \item[\textbf{14.}] $(-\infty,-1)\cup(1,+\infty)$。机器人轨迹为抛物线 $y^2=4x$；直线 $y=k(x+1)$ 与它无公共点，判别式小于零给出 $k^2>1$。
  \item[\textbf{15.}] $-\dfrac32$。由 $f(-x)=f(x)$ 得 $\ln(\mathrm e^{3x}+1)-3x-ax=\ln(\mathrm e^{3x}+1)+ax$，故 $a=-\dfrac32$。
\end{itemize}

\gkpart{解答题}
\begin{itemize}[leftmargin=2.6em]
  \item[\textbf{16.}] (1) $a_n=n$。$a_1=S_1=1$，$n\geqslant2$ 时 $a_n=S_n-S_{n-1}=n$，对 $n=1$ 也成立；(2) $2^{2n+1}+n-2$。$b_n=2^n+(-1)^nn$，等比部分求和为 $2^{2n+1}-2$，交错部分求和为 $n$
  \item[\textbf{17.}] (1) 甲组平均数 $\dfrac23$、方差 $\dfrac29$；乙组平均数 $\dfrac35$、方差 $\dfrac6{25}$；甲组平均数较高且方差较小，研发水平较高。$15$ 次结果中甲成功 $10$ 次、乙成功 $9$ 次；(2) $\dfrac7{15}$。以频率估计概率，$P=\dfrac23\cdot\dfrac25+\dfrac13\cdot\dfrac35=\dfrac7{15}$
  \item[\textbf{18.}] (1) 见解析版；(2) $\dfrac34$。$BC\parallel AD$，建系求得 $\cos\langle\overrightarrow{AD},\overrightarrow{OD}\rangle=\dfrac34$
  \item[\textbf{19.}] (1) $\dfrac{\sqrt{21}}7$。$\angle CDE=120^\circ$，由余弦定理得 $DC=2$，再由正弦定理得 $\sin\angle CED=\dfrac{2\sin120^\circ}{\sqrt7}=\dfrac{\sqrt{21}}7$；(2) $4\sqrt7$。$\cos\angle AEB=\cos\left(\angle AEC-60^\circ\right)=\dfrac1{2\sqrt7}$，在 Rt$\triangle ABE$ 中 $BE=\dfrac{EA}{\cos\angle AEB}=4\sqrt7$
  \item[\textbf{20.}] (1) $C_1:x^2-\dfrac{y^2}3=1$；$C_2:\dfrac{y^2}3+\dfrac{x^2}2=1$。正方形对角线相等且面积为 $2$ 给出 $a_1=c_2=1$，再代入 $P$ 定出 $b_1^2=3$，$b_2^2=2$；(2) 不存在。条件等价于 $\overrightarrow{OA}\perp\overrightarrow{OB}$；与 $C_2$ 相切给出 $m^2=2k^2+3$，而 $\overrightarrow{OA}\cdot\overrightarrow{OB}=0$ 要求 $2m^2=3(1+k^2)$，联立得 $k^2=-3$，无解
  \item[\textbf{21.}] (1) 单调递减区间为 $(2k\pi,(2k+1)\pi)$，单调递增区间为 $((2k+1)\pi,(2k+2)\pi)$，$k\in\mathbf N$。$f'(x)=-x\sin x$；(2) 见解析版
\end{itemize}
